\section*{Hoofdstuk \ref{ch:b3:organometallic-bonding} --- Organometaalchemie: binding en liganden}

\begin{solution}{exo:b3:organometallic-bonding:1}
$\ce{Fe(CO)5}$: $8 + 10 = 18$. $\ce{Mn2(CO)10}$: $7 + 10 + 1 = 18$ per Mn. $\ce{CpMo(CO)3H}$: $6 + 5 + 6 + 1 =
18$. $\ce{[PtCl3(C2H4)]-}$: $10 + 3 + 2 + 1 = 16$. $\ce{Cr($\eta^6$-C6H6)(CO)3}$: $6 + 6 + 6 = 18$. $\ce{Cp2ZrCl2}$:
$4 + 10 + 2 = 16$.
\end{solution}

\begin{solution}{exo:b3:organometallic-bonding:2}
Fe(0) $d^8$; Mn(0) $d^7$; Mo(II) $d^4$; Pt(II) $d^8$; Cr(0) $d^6$; Zr(IV) $d^0$.
\end{solution}

\begin{solution}{exo:b3:organometallic-bonding:3}
$\ce{[Mn(CO)6]+} > \ce{Cr(CO)6} > \ce{[V(CO)6]-}$: hoe elektronenrijker het metaal, hoe meer het
terugdoneert in $\pi^*$ en hoe zwakker de bindingen C--O.
\end{solution}

\begin{solution}{exo:b3:organometallic-bonding:4}
Ethaan ($\ce{Mn(CO)5}$ en $\ce{CpFe(CO)2}$ zijn \omterm{def:b3:organometallic-bonding:isolobal}{isolobaal} met \ce{CH3}); opnieuw ethaan; tetraëdraan
$\ce{(CH)4}$ waarin drie CH vervangen zijn door $\ce{Co(CO)3}$.
\end{solution}

\begin{solution}{exo:b3:organometallic-bonding:5}
$fac$ ($C_{3v}$): twee (A$_1$ + E). $mer$ ($C_{2v}$): drie (2A$_1$ + B$_1$).
\end{solution}

\begin{solution}{exo:b3:organometallic-bonding:6}
Het omvangrijke fosfine verdringt het metaal: dissociatie heft de
spanning op en verloopt dus sneller, en laaggecoördineerde intermediairen
worden begunstigd.
\end{solution}

\begin{solution}{exo:b3:organometallic-bonding:7}
Het chroomcomplex heeft een \omterm{def:b3:organometallic-bonding:carbene}{fischercarbeen} (substituent OMe, Cr(0), liganden CO): een
amine valt het elektrofiele carbeenkoolstofatoom aan en vervangt OMe. Het
tantaalalkylideen is een \omterm{def:b3:organometallic-bonding:carbene}{schrockcarbeen}: zijn nucleofiele koolstofatoom
valt het carbonyl van een
keton aan, wat een alkeen en een eenheid Ta=O geeft, zoals bij een wittigreactie.
\end{solution}

\begin{solution}{exo:b3:organometallic-bonding:8}
Verspringend: geen $\delta$-overlap, de twee $\delta$-elektronen zijn niet-bindend: bindingsorde 3. $\ce{[Re2Cl8]^{4-}}$:
$\sigma^2\pi^4\delta^2\delta^{\ast2}$, bindingsorde 3.
\end{solution}

\begin{solution}{exo:b3:organometallic-bonding:9}
Een korte afstand M$\cdots$H (diffractie, liefst met neutronen); een laag
golfgetal van de strekvibratie C--H; een NMR-signaal bij hoog veld met
een verkleinde koppeling C--H over één binding.
\end{solution}

\begin{solution}{exo:b3:organometallic-bonding:10}
16: vlak vierkante $d^8$ ($\ce{[PtCl4]^{2-}}$, de katalysator van Wilkinson), waarvan de orbitaal voor het 18e elektron
te hoog ligt; metallocenen van vroege metalen zoals $\ce{Cp2ZrCl2}$, te vol voor meer
liganden. 17: $\ce{V(CO)6}$, dat om sterische redenen niet kan dimeriseren. 19: cobaltoceen,
waarvan het extra elektron in een antibindende orbitaal zit (en
gemakkelijk verloren gaat).
\end{solution}

\begin{solution}{exo:b3:organometallic-bonding:11}
CO is een sterke $\pi$-acceptor: het verlaagt het $t_{2g}$-niveau en maakt $\Delta_{\mathrm o}$ groot, ver boven de
paringsenergie: $t_{2g}^6$, diamagnetisch. De $d$--$d$-overgangen, bij $\Delta_{\mathrm o}$, vallen in het ultraviolet,
en er wordt geen zichtbaar licht geabsorbeerd.
\end{solution}

\begin{solution}{exo:b3:organometallic-bonding:12}
Donatie vanuit $\pi$ en terugdonatie in $\pi^*$ verlagen allebei de bindingsorde C--C. In het zout
van Zeise (Pt(II), matige terugdonatie) wordt de binding iets langer; een
laagvalent, elektronenrijk metaal doneert sterk terug en het complex
nadert de limiet van het \omterm{def:b3:organometallic-bonding:dcd}{metallacyclopropaan}, met een binding C--C dicht
bij een enkelvoudige binding.
\end{solution}

\begin{solution}{pb:b3:organometallic-bonding:1}
\textbf{1.} $\ce{Fe^{2+}}$ ($d^6$) $+ 2 \times 6 = 18$.
\textbf{2.} $8 + 2 \times 5 = 18$.
\textbf{3.} $9 + 10 = 19$.
\textbf{4.} $10 + 10 = 20$.
\textbf{5.} $18 - 1 = 17$.
\textbf{6.} Fe(II) $d^6$; Co(II) $d^7$; Ni(II) $d^8$; Fe(III) $d^5$.
\textbf{7.} $a_{1g}$ ($d_{z^2}$), $e_{1g}$ ($d_{xz}$, $d_{yz}$), $e_{2g}$ ($d_{xy}$, $d_{x^2-y^2}$).
\textbf{8.} Het paar $e_{1g}$ overlapt sterk met de gevulde ringcombinatie $e_{1g}$: bindende
orbitalen vooral op de ringen, antibindende $e_{1g}^\ast$ vooral op het metaal, hoog in
energie.
\textbf{9.} $e_{2g} \approx a_{1g} < e_{1g}^\ast$.
\textbf{10.} $e_{2g}^4a_{1g}^2$: geen ongepaard elektron.
\textbf{11.} Eén elektron in $e_{1g}^\ast$: één ongepaard elektron, $1.73\,\mu_B$.
\textbf{12.} Twee elektronen in de twee $e_{1g}^\ast$-orbitalen, parallel: twee ongepaarde, $\sqrt{8} =
2.83\,\mu_B$.
\textbf{13.} $e_{2g}^3a_{1g}^2$: één ongepaard elektron. Een gat in $e_{2g}^3$ geeft een orbitaal ontaarde
toestand (E) met een \omterm{def:b3:complex-spectra-magnetism:moment}{orbitale bijdrage}.
\textbf{14.} Zijn 19e elektron zit in een antibindende orbitaal en wordt
gemakkelijk verwijderd, wat het cobaltoceniumkation met 18 elektronen
geeft.
\textbf{15.} Het verwijderde elektron komt uit een nagenoeg niet-bindende
orbitaal: de structuur verandert nauwelijks (kleine \omterm{def:b3:complex-mechanisms:reorganisation}{reorganisatie-energie}, \cref{ch:b3:complex-mechanisms}).
\textbf{16.} Het koppel is in de meeste oplosmiddelen snel en reversibel,
de verbinding is oplosbaar en stabiel, en haar potentiaal hangt weinig af
van het oplosmiddel omdat de lading over een groot ion verdeeld is.
\textbf{17.} Met 20 elektronen, waarvan twee antibindend, reageert het door
ze kwijt te raken: het bindt liganden onder verglijden of verlies van een
ring, of het wordt geoxideerd.
\textbf{18.} $\ce{CpFe(CO)2}$: $8 + 5 + 4 = 17$ elektronen, één te weinig voor 18, één grensorbitaal: \omterm{def:b3:organometallic-bonding:isolobal}{isolobaal}
met \ce{CH3}. Het dimeer $\ce{[CpFe(CO)2]2}$ is ``ethaan'', met een binding Fe--Fe (in de praktijk overbruggen twee
CO-liganden haar).
\textbf{19.} $\ce{CpFe(CO)2-Mn(CO)5}$, met een binding Fe--Mn.
\textbf{20.} $\ce{PPh3}$ is een betere donor en een zwakkere $\pi$-acceptor dan CO: het metaal doneert
meer terug naar de overige CO, waarvan de strekfrequenties dalen.
\textbf{21.} Twee (A$_1$ + E).
\textbf{22.} Het verglijden van de ring naar $\eta^3$ maakt twee elektronen aan coördinatie vrij: een binnenkomend
ligand kan associatief binden zonder dat het metaal boven 18 elektronen
uitkomt.
\textbf{23.} \textbf{$\mu_{\text{spin-only}}(\text{nikkeloceen}) = \sqrt{2 \times 4} = 2.83\,\mu_B$.}
\end{solution}
